\documentclass[12pt]{article}

\usepackage{amsmath,amssymb,amsthm,mathtools}
\usepackage[margin=1in]{geometry}
\usepackage{enumitem}
\usepackage{xurl}
\usepackage[hidelinks]{hyperref}

\theoremstyle{plain}
\newtheorem{theorem}{Theorem}
\newtheorem{proposition}[theorem]{Proposition}
\newtheorem{lemma}[theorem]{Lemma}
\newtheorem{corollary}[theorem]{Corollary}
\theoremstyle{definition}
\newtheorem{remark}[theorem]{Remark}
\theoremstyle{plain}

\newcommand{\val}[2]{[#1]_{#2}}
\newcommand{\NpD}{\mathcal{N}_{\mathrm{pD}}}
\newcommand{\DP}{\mathcal{D}_{\mathrm{P}}}
\newcommand{\Pset}{%
  \bigl(\tfrac{1}{4}\cdot\mathcal{N}_{\mathrm{pD},\geq 12}\bigr)%
}
\newcommand{\seqnum}[1]{\href{https://oeis.org/#1}{\rm \underline{#1}}}

\title{Additive Representations by Primitive Dyck Numbers:\\ Exceptional Sets and Optimal Orders}
\author{Takayuki Kuriyama}
\date{}

\begin{document}

\maketitle

\begin{abstract}
Let $\NpD$ be the set consisting of $0$ and the integers represented in
binary by primitive Dyck words.  For a positive even integer $N$, let $r(N)$
be the least number of positive elements of $\NpD$ whose sum is $N$; the
associated sequence $a(n)=r(2n)$ is OEIS sequence \seqnum{A395858}.  We
determine the complete exceptional set for six-summand representability:
$46$ alone satisfies $r(46)=8$, the ten integers
$34,44,98,154,198,202,206,838,842,846$ satisfy $r(N)=7$, and every other
positive even integer satisfies $r(N)\leq6$.  Thus $848$ is the sharp even
threshold beyond which six summands always suffice.  We also determine the
exact relative asymptotic order.  For every $k\geq2$, one has
$r(10\cdot4^{k+1}-6)=6$, so six summands are necessary infinitely often and the asymptotic order is
exactly $6$.  The proof combines a regular sublanguage, finite interval
certificates, a general digit-lifting mechanism that propagates finite
sumset certificates to infinite tails, and a base-$4$ characterization of
the halved primitive family by two-coloured Motzkin words.
As a further consequence, that halved family is an asymptotic additive basis
of order exactly $5$.
\end{abstract}

\noindent\textbf{Keywords:} primitive Dyck number, additive
representation, additive basis, asymptotic order, exceptional set,
two-coloured Motzkin word, digital self-similarity, interval filling.

\noindent\textbf{2020 Mathematics Subject Classification:} Primary 11B13;
Secondary 05A15, 68Q45.

\section{Introduction}\label{sec:introduction}

For example, the balanced parenthesis string $()(())$ becomes the word
$101100$ when an opening parenthesis is encoded by $1$ and a closing
parenthesis by $0$.  More generally, every balanced parenthesis string can
be written over the alphabet $\{1,0\}$ in this way\footnote{The set of all
words over this alphabet is denoted by $\{0,1\}^*$; it is a regular
language.} and hence read as a binary integer.  How additively rich is the
set of integers obtained in this way?  We take up a natural variant of this
question.

A \emph{Dyck word} is a word $w\in\{1,0\}^*$ with equally many $1$'s
and $0$'s in which every prefix has at least as many $1$'s as $0$'s.\footnote{The
Dyck language is not regular, but it is context-free.}  For example, the
unique factorization of the Dyck word $101100$ into primitive Dyck words is
$101100=10\cdot1100$.  Here and below, $\cdot$ denotes concatenation of
words.  A nonempty Dyck word is \emph{primitive} if none of
its proper nonempty prefixes is balanced.  Equivalently, if $1$ is viewed as
an up-step and $0$ as a down-step, the associated lattice path returns to
height zero for the first time only at its endpoint.\footnote{Such paths are
also commonly called \emph{irreducible} or \emph{indecomposable} Dyck paths.}
Let $\DP$ denote the language of nonempty primitive Dyck words; in
particular, the empty word $\lambda$ is not included in $\DP$.  Ordered by length and then
lexicographically, its first eighteen elements are
\[
\begin{gathered}
10,\ 1100,\ 110100,\ 111000,\ 11010100,\ 11011000,\\
11100100,\ 11101000,\ 11110000,\ 1101010100,\ 1101011000,\\
1101100100,\ 1101101000,\ 1101110000,\ 1110010100,\\
1110011000,\ 1110100100,\ 1110101000.
\end{gathered}
\]

Binary words of even length may also be read blockwise in base $4$.
For example, at the level of digit strings we write
\[
 (11\cdot10\cdot01\cdot01\cdot00)_{(2)}=32110_{(4)},
\]
where $\cdot$ denotes concatenation.  More generally, consecutive binary
pairs are identified with base-$4$ digits through
\[
 11_{(2)}\longleftrightarrow3_{(4)},\qquad
 10_{(2)}\longleftrightarrow2_{(4)},\qquad
 01_{(2)}\longleftrightarrow1_{(4)},\qquad
 00_{(2)}\longleftrightarrow0_{(4)}.
\]
The subscripts $(2)$ and $(4)$ are metanotational base labels and will be
omitted whenever the base is clear from context.

For a word $w$ over the base-$m$ digits, let $\val{w}{m}$ denote the
integer it represents in base $m$.  For convenience in the additive
representations considered in this paper, we adjoin $0$ and set
\[
 \NpD=\{0\}\cup\{\val{w}{2}:w\in\DP\}
 \subseteq\mathbb Z_{\geq0}.
\]
The convention $0\in\NpD$ is inessential to the conclusions below.  The
first eighteen positive elements are
\[
\begin{gathered}
2,12,52,56,212,216,228,232,240,\\
852,856,868,872,880,916,920,932,936;
\end{gathered}
\]
this is OEIS sequence \seqnum{A057547} \cite{OEIS-A057547}.  Every positive
member of $\NpD$ is even, since every nonempty Dyck word ends in $0$.
Thus only even target integers can be represented by sums from $\NpD$.

We shall repeatedly use the following elementary mod-$4$ observation.

\begin{lemma}\label{lem:mod4}
The number $2$ is the only positive primitive Dyck number congruent to
$2$ modulo $4$.  Every other positive primitive Dyck number is divisible
by $4$.
\end{lemma}

\begin{proof}
Let $w\in\DP$, and let $|w|$ denote its length.  Since $w$ is balanced and
nonempty, its last letter is $0$.  If $|w|=2$, then necessarily $w=10$ and
$[w]_2=2$.  Suppose that $|w|\geq4$ and write
\[
 w=ub0,\qquad b\in\{0,1\}.
\]
If $b=1$, then the suffix $10$ has net height zero, so the proper prefix
$u$ is already balanced, contradicting the primitivity of $w$.  Hence
$b=0$.  Thus $w$ ends in $00$, and therefore $4\mid[w]_2$.
\end{proof}

Bell, Lidbetter, and Shallit studied the larger set obtained from
\emph{all} Dyck words, the \emph{totally balanced numbers} (OEIS sequence
\seqnum{A014486} \cite{OEIS-A014486}).  Using regular underapproximations and automata,
they proved that every even integer except
\[
        8,\ 18,\ 28,\ 38,\ 40,\ 82,\ 166
\]
is a sum of at most three totally balanced numbers, while two summands
do not suffice asymptotically \cite[Thm.\ 14]{BLS}.

The primitive restriction is not a cosmetic variant of that result.
Every nonempty Dyck word factors uniquely as a concatenation of
primitive Dyck words, whereas here each summand must consist of a single
factor.  Consequently, a representation by unrestricted Dyck numbers
does not in general yield a representation by the same number of
primitive Dyck numbers.

The additive representations considered here are far from unique.  For
example,
\[
 2026=2+240+852+932=2+216+872+936.
\]
The seven distinct integers appearing in these two representations all
have primitive Dyck binary expansions.

Digit-defined additive bases have also been studied by automata-theoretic
methods.  Bell, Hare, and Shallit characterized automatic sets that form
additive bases and gave an effective procedure for determining their
orders \cite{BHS}.  Related decision procedures have yielded sharp
additive results for binary palindromes \cite{RSS} and binary squares
\cite{MNRS}.  The language of primitive Dyck words is not regular, so
those general finite-automaton methods do not directly settle the
present problem.

Our proof differs at the structural point where regular approximation
alone ceases to suffice.  A regular sublanguage settles the congruence
class $0$ modulo $4$, while the class $2$ modulo $4$ is handled by a
self-similar closure property of the full primitive family after division
by $4$.  The same family admits a base-$4$ description by two-coloured
Motzkin words.  Its self-similarity has two complementary effects: it
propagates five-summand intervals, but its generation gaps also create an
infinite sequence of obstructions to four-summand representations.  These
facts yield the complete exceptional set, the optimal uniform bound, the
sharp eventual threshold, and the exact asymptotic order.

Within the Dyck words of length $2n$, the proportion that are primitive
tends to $1/4$ as $n\to\infty$.  Indeed, there are $C_n$ Dyck words of
length $2n$, and every primitive one is uniquely of the form $1u0$ with
$u$ a Dyck word of length $2n-2$; so there are $C_{n-1}$ primitive words
of length $2n$.  From the Catalan formula
$C_n=\frac1{n+1}\binom{2n}{n}$\footnote{The context-free grammar
$S\to\lambda\mid 1S0S$ gives the unique decomposition $1u0v$ of every
nonempty Dyck word, where $u$ and $v$ are Dyck words.  Hence $C_0=1$ and
$C_n=\sum_{k=0}^{n-1}C_kC_{n-1-k}$.  If
$B_n=\frac{1}{n+1}\binom{2n}{n}$, then the binomial convolution identity gives
$B_0=1$ and $B_n=\sum_{k=0}^{n-1}B_kB_{n-1-k}$.  Therefore induction yields
$C_n=B_n$.},
\[
 \frac{C_{n-1}}{C_n}=\frac{n+1}{2(2n-1)}\longrightarrow\frac14 .
\]
Its arithmetic behaviour, however, acquires a genuinely hierarchical
feature.  Although all sufficiently large even integers need only three
unrestricted Dyck summands, restricting the summands to the primitive family
produces a small but nontrivially shaped exceptional set.  We show that
precisely eleven positive even integers fail to be sums of at most six
primitive Dyck numbers.  Ten require seven summands, while $46$ alone
requires eight.  All eleven exceptions lie below $848$.

The eight-summand phenomenon is local and is caused by the fact that the only
positive primitive Dyck numbers below $48$ are $2$ and $12$.  By contrast,
the sharpness of the eventual six-summand bound is structural.  A regular
subfamily handles multiples of $4$, while the full primitive family after
division by $4$ handles the remaining congruence class; moreover, the
explicit sequence $(N_k)_{k\geq2}$ constructed below tends to infinity and
satisfies $r(N_k)=6$.

For a positive even integer $N$, let $r(N)$ be the least number of
positive primitive Dyck numbers whose sum is $N$, and define the central
integer sequence of this paper by
\[
        a(n)=r(2n),\qquad n\geq1.
\]
This sequence has been registered in the OEIS as
\seqnum{A395858} \cite{OEIS-A395858}.  Its first fifty terms are
\[
\begin{gathered}
1,2,3,4,5,1,2,3,4,5,6,2,3,4,5,6,7,3,4,5,6,7,8,4,5,\\
1,2,1,2,3,4,2,3,2,3,4,5,3,4,3,4,5,6,4,5,4,5,6,7,5.
\end{gathered}
\]
The generating program is
\href{https://github.com/growupkuriyama-hub/lean_cfg_project/blob/3525cfb/LeanCfgProject/PrimDyck/A395858.py}{\path{A395858.py}},
and the actual output of its first $50{,}000$ terms is available as
\href{https://github.com/growupkuriyama-hub/lean_cfg_project/blob/3525cfb/LeanCfgProject/PrimDyck/b395858.txt}{\path{b395858.txt}}.

\begin{theorem}\label{thm:classification}
The values of $r(N)$ exceeding six are exactly as follows:
\[
\begin{aligned}
 r(46)&=8,\\[4pt]
 r(N)=7\quad&\Longleftrightarrow\quad
 N\in\{34,44,98,154,198,202,206,838,842,846\}.
\end{aligned}
\]
Thus every other positive even integer satisfies $r(N)\leq6$.  In other
words, every even integer $N\geq848$ is a sum of at most six primitive
Dyck numbers.
\end{theorem}

Equivalently, for the OEIS sequence \seqnum{A395858}
\cite{OEIS-A395858}, defined by $a(n)=r(2n)$,
\[
        \max_{n\geq1}a(n)=8,
        \qquad
        a(n)=8\quad\Longleftrightarrow\quad n=23,
\]
and
\[
 a(n)=7
 \quad\Longleftrightarrow\quad
 n\in\{17,22,49,77,99,101,103,419,421,423\}.
\]
Thus the main theorem gives a complete determination of all terms of
$a(n)$ exceeding six, rather than only an eventual estimate.


A casual inspection of the computed terms may suggest that occurrences of
the value $6$ become rarer and eventually disappear.  This is not the case:
there is an explicit sequence of increasingly remote obstructions $N_k$, and
the eventual upper bound six is best possible.

\begin{theorem}\label{thm:asymptotic-order}
Let
\[
 h_{\mathrm{asym}}
 =\min\{h:\text{every sufficiently large even }N\text{ satisfies }r(N)\leq h\}.
\]
Then $h_{\mathrm{asym}}=6$.  More precisely, for every integer $k\geq2$,
\[
 N_k=10\cdot4^{k+1}-6
\]
satisfies $r(N_k)=6$.
\end{theorem}

Thus the finite exceptional set above coexists with infinitely many large
integers for which the six-summand upper bound is attained.

The exceptional value $46$ already explains why the uniform bound can
never be smaller than eight.

\begin{proposition}\label{prop:46}
The integer $46$ requires exactly eight positive primitive Dyck numbers:
it is a sum of eight, but it is not a sum of at most seven.
\end{proposition}

\begin{proof}
Below $48$ the only positive primitive Dyck numbers are $2$ and $12$.
Indeed, the primitive words of length two and four are $10$ and $1100$,
whereas every primitive word of length at least six begins with $11$ and
therefore represents an integer at least $[110000]_2=48$.  Thus any
representation of $46$ has the form
\[
        46=12a+2b,\qquad a,b\geq0.
\]
Equivalently, $6a+b=23$.  Since $6\cdot4>23$, we have $a\leq3$, and
hence the number of summands satisfies
\[
        a+b=23-5a\geq8.
\]
Equality is attained by
\[
        46=12+12+12+2+2+2+2+2.
\]
\end{proof}


The sumset notation used below is standard in additive-basis theory
\cite[Chap.\ 1]{Nathanson}.  For a nonnegative integer $k$, let $kX$ denote the
set of sums of $k$ not necessarily distinct members of $X$, with
$0X=\{0\}$, and put
\[
 F_k(X)=\bigcup_{j=0}^{k}jX.
\]
Thus $F_k(X)$ is the set of sums of at most $k$ elements of $X$.  This
notation should not be confused with dilation: in general,
$jX\neq\{jx:x\in X\}$.  For a scalar $c$, we distinguish the dilation by
writing $c\cdot X=\{cx:x\in X\}$.

The mod-$4$ property motivates the central halved family
\begin{equation}\label{eq:P-def}
 \Pset=\{v/4:v\in\NpD,\ v\geq12\}.
\end{equation}
Its first elements are
\[
 3,13,14,53,54,57,58,60,213,214,217,\ldots.
\]
Then
\begin{equation}\label{eq:NpD-decomp}
 \NpD=\{0,2\}\cup4\Pset.
\end{equation}


\section{The base-\texorpdfstring{$4$}{4} structure of
\texorpdfstring{$\Pset$}{P}}\label{sec:base4}

We now describe the base-$4$ structure of $\Pset$, which is exposed by
the division by $4$ in its definition.  Assign weights to the base-$4$ digits by
\[
  \sigma(3)=1,\qquad \sigma(0)=-1,\qquad
  \sigma(1)=\sigma(2)=0.
\]
Extend $\sigma$ additively to words by
\[
 \sigma(c_0c_1\cdots c_{k-1})
 :=\sum_{i=0}^{k-1}\sigma(c_i),
 \qquad
 \sigma(\lambda):=0.
\]
A word $w=c_0\cdots c_{k-1}\in\{0,1,2,3\}^{k}$ is called a
\emph{two-coloured Motzkin word}\footnote{The two zero-weight digits $1$
and $2$ are regarded as distinct colours.} if every prefix has nonnegative
weight and $\sigma(w)=0$.  The running prefix weight is allowed to return to
$0$ more than once.  Let $\mathsf{Motz}$ denote the set of all such words.
Unlike $\DP$, the set $\mathsf{Motz}$ contains the empty word $\lambda$.
Write $[w]_4$ for the value of $w$ as a base-$4$ numeral, and set
$[\lambda]_4=0$.  For a positive integer $n$, let
$\operatorname{bin}(n)$ and $\operatorname{quad}(n)$ denote its standard
binary and base-$4$ expansions, respectively.

\begin{lemma}[Base-$4$ characterization]\label{lem:base4}
A positive integer $p$ lies in $\Pset$ if and only if there exists
$w\in\mathsf{Motz}$ such that $\operatorname{quad}(p)=3w$.
\end{lemma}

\begin{proof}
We give the details explicitly.  Suppose first that $p\in\Pset$.  By the definition of $\Pset$,
\[
 4p\in\NpD,\qquad 4p\geq12.
\]
Hence the binary expansion
\[
 B:=\operatorname{bin}(4p)=b_1b_2\cdots b_{2L},
 \qquad b_i\in\{0,1\},
\]
is a primitive Dyck word satisfying $|B|=2L\geq4$ and $B\in\DP$.
Define its height function by
\[
 h_B(t):=\sum_{i=1}^{t}(2b_i-1)
 \qquad(0\leq t\leq2L).
\]
The condition $B\in\DP$ is equivalently expressed by
\begin{equation}\label{eq:primitive-height}
 h_B(0)=h_B(2L)=0,
 \qquad
 h_B(t)>0\quad(1\leq t<2L).
\end{equation}

Write the base-$4$ expansion of $4p$ as
\[
 \operatorname{quad}(4p)=q_1q_2\cdots q_L,
 \qquad q_j\in\{0,1,2,3\}.
\]
Grouping the binary digits from the left into consecutive pairs gives
\[
 q_j=2b_{2j-1}+b_{2j}
 \qquad(1\leq j\leq L).
\]
The correspondences
\[
 3_{(4)}\longleftrightarrow11_{(2)},\qquad
 2_{(4)}\longleftrightarrow10_{(2)},\qquad
 1_{(4)}\longleftrightarrow01_{(2)},\qquad
 0_{(4)}\longleftrightarrow00_{(2)}
\]
show that the net height change across the $j$th two-bit block is twice
the weight of the corresponding base-$4$ digit:
\[
 h_B(2j)-h_B(2j-2)=2\sigma(q_j).
\]
Summing these identities from $1$ to $j$ yields
\begin{equation}\label{eq:height-sigma}
 h_B(2j)=2\sum_{i=1}^{j}\sigma(q_i)
 \qquad(0\leq j\leq L),
\end{equation}
where the sum is empty when $j=0$.

We next determine the first and last blocks.  Since $B$ is a Dyck word,
$b_1=1$.  If $b_2=0$, then $h_B(2)=0$, contradicting
$2<2L$ and \eqref{eq:primitive-height}.  Hence
\[
 b_1b_2=11_{(2)},
 \qquad
 q_1=2b_1+b_2=3.
\]
Similarly, $b_{2L}=0$.  If $b_{2L-1}=1$, then the last block is $10$, so
\[
 h_B(2L)=h_B(2L-2).
\]
Since $h_B(2L)=0$, this would give $h_B(2L-2)=0$, again contradicting
\eqref{eq:primitive-height}.  Thus
\[
 b_{2L-1}b_{2L}=00_{(2)},
 \qquad
 q_L=0.
\]

It remains to verify the weight conditions defining $\mathsf{Motz}$.
For $1\leq j\leq L-1$, one has $2j<2L$, so
$h_B(2j)>0$ by \eqref{eq:primitive-height}.  Since $h_B(2j)$ is a positive
even integer, equation~\eqref{eq:height-sigma} gives
\begin{equation}\label{eq:prefix-height}
 \sum_{i=1}^{j}\sigma(q_i)
 =\frac{h_B(2j)}{2}
 \geq1
 \qquad(1\leq j\leq L-1).
\end{equation}
Dividing $4p$ by $4$ deletes the terminal base-$4$ digit $q_L=0$.
Since $q_1=3$, we obtain
\[
 \operatorname{quad}(p)=q_1q_2\cdots q_{L-1}=3u,
 \qquad
 u:=q_2q_3\cdots q_{L-1}.
\]

It remains to show that $u\in\mathsf{Motz}$.  Consider a prefix of $u$ of
length $r$, where $0\leq r\leq L-2$.  For $r=0$ its weight is the empty sum
$0$.  For $1\leq r\leq L-2$, applying \eqref{eq:prefix-height} with
$j=r+1$ gives
\begin{equation}\label{eq:motzkin-prefix-nonnegative}
 \sum_{i=2}^{r+1}\sigma(q_i)
 =\sum_{i=1}^{r+1}\sigma(q_i)-\sigma(q_1)
 \geq1-1=0.
\end{equation}
Moreover, applying \eqref{eq:height-sigma} with $j=L$ gives
\begin{equation}\label{eq:total-weight-zero}
 \sum_{i=1}^{L}\sigma(q_i)
 =\frac{h_B(2L)}{2}
 =0.
\end{equation}
Since $\sigma(q_1)=\sigma(3)=1$ and
$\sigma(q_L)=\sigma(0)=-1$, equation~\eqref{eq:total-weight-zero} gives
\[
 \sum_{i=2}^{L-1}\sigma(q_i)
 =\sum_{i=1}^{L}\sigma(q_i)
  -\bigl(\sigma(q_1)+\sigma(q_L)\bigr)
 =0.
\]
Equation~\eqref{eq:motzkin-prefix-nonnegative} shows that every prefix
of $u$ has nonnegative weight, while \eqref{eq:total-weight-zero} and the
endpoint weights $\sigma(q_1)=1$ and $\sigma(q_L)=-1$ show that the total
weight of $u$ is $0$.  Hence $u\in\mathsf{Motz}$, proving the forward
implication.

Conversely, suppose that there exists $w\in\mathsf{Motz}$ such that
$\operatorname{quad}(p)=3w$.  Write
\[
 \operatorname{quad}(p)=q_1q_2\cdots q_m,
 \qquad
 q_1=3,
 \qquad
 w=q_2q_3\cdots q_m.
\]
The condition $w\in\mathsf{Motz}$ is expressed explicitly by
\begin{equation}\label{eq:reverse-motzkin-prefix}
 \sum_{i=2}^{j}\sigma(q_i)\geq0
 \qquad(2\leq j\leq m)
\end{equation}
and
\begin{equation}\label{eq:reverse-motzkin-total}
 \sum_{i=2}^{m}\sigma(q_i)=0.
\end{equation}
When $m=1$, both sums are understood as empty and the prefix condition is
vacuous.

Define the monoid morphism
\[
 \phi:\{0,1,2,3\}^*\longrightarrow\{0,1\}^*
\]
by
\[
 \phi(3)=11,\qquad
 \phi(2)=10,\qquad
 \phi(1)=01,\qquad
 \phi(0)=00
\]
on letters.  Put $q_{m+1}:=0$ and set
\[
 B:=\phi(q_1q_2\cdots q_mq_{m+1})
   =b_1b_2\cdots b_{2m+2}.
\]
The final two bits satisfy
\[
 b_{2m+1}b_{2m+2}=00_{(2)}.
\]
Because each base-$4$ digit is replaced by its two-bit binary block,
\begin{equation}\label{eq:block-value}
 [B]_2=[q_1q_2\cdots q_m0]_4=4p.
\end{equation}
Define
\[
 h_B(t):=\sum_{i=1}^{t}(2b_i-1)
 \qquad(0\leq t\leq2m+2).
\]
The same block calculation as above gives
\begin{equation}\label{eq:reverse-height-sigma}
 h_B(2j)=2\sum_{i=1}^{j}\sigma(q_i)
 \qquad(0\leq j\leq m+1).
\end{equation}

We first check the block boundaries.  For $1\leq j\leq m$,
\[
 \sum_{i=1}^{j}\sigma(q_i)
 =\sigma(q_1)+\sum_{i=2}^{j}\sigma(q_i)
 =1+\sum_{i=2}^{j}\sigma(q_i)
 \geq1,
\]
where the second sum is empty when $j=1$.  Hence
\begin{equation}\label{eq:positive-even-height}
 h_B(2j)\geq2
 \qquad(1\leq j\leq m).
\end{equation}
On the other hand, $\sigma(q_{m+1})=\sigma(0)=-1$ and
\eqref{eq:reverse-motzkin-total} imply
\[
 \sum_{i=1}^{m+1}\sigma(q_i)
 =\sigma(q_1)+\sum_{i=2}^{m}\sigma(q_i)+\sigma(q_{m+1})
 =1+0-1=0.
\]
Thus
\begin{equation}\label{eq:terminal-height}
 h_B(2m+2)=0.
\end{equation}

It remains to check the positions inside the blocks, namely the odd positions
$2j-1$ for $1\leq j\leq m+1$.  For $j=1$, since $q_1=3$, the first block
is $11$, and hence $h_B(1)=1$.

Now let $2\leq j\leq m$.  If $q_j\in\{2,3\}$, its block is $10$ or $11$,
so its first bit is $1$.  By \eqref{eq:positive-even-height}, applied with
$j-1$ in place of $j$,
\[
 h_B(2j-1)=h_B(2j-2)+1\geq3.
\]
If $q_j=1$, its block is $01$, and the same shifted application of
\eqref{eq:positive-even-height} gives
\[
 h_B(2j-1)=h_B(2j-2)-1\geq1.
\]
Finally suppose that $q_j=0$.  By \eqref{eq:reverse-motzkin-prefix},
\[
 0\leq\sum_{i=2}^{j}\sigma(q_i)
 =\sum_{i=2}^{j-1}\sigma(q_i)+\sigma(q_j)
 =\sum_{i=2}^{j-1}\sigma(q_i)-1,
\]
so
\[
 \sum_{i=2}^{j-1}\sigma(q_i)\geq1.
\]
Applying \eqref{eq:reverse-height-sigma} with $j-1$ in place of $j$,
and using $\sigma(q_1)=1$, gives
\[
\begin{aligned}
 h_B(2j-2)
 &=2\sum_{i=1}^{j-1}\sigma(q_i)\\
 &=2\left(1+\sum_{i=2}^{j-1}\sigma(q_i)\right)
 \geq4.
\end{aligned}
\]
and hence
\[
 h_B(2j-1)=h_B(2j-2)-1\geq3.
\]

For the appended terminal digit $q_{m+1}=0$, equations
\eqref{eq:reverse-height-sigma} and
\eqref{eq:reverse-motzkin-total} give
\[
\begin{aligned}
 h_B(2m)
 &=2\left(\sigma(q_1)+\sum_{i=2}^{m}\sigma(q_i)\right)\\
 &=2\left(1+\sum_{i=2}^{m}\sigma(q_i)\right)
 =2.
\end{aligned}
\]
Thus the terminal block $00$ passes through height
\[
 h_B(2m+1)=1>0
\]
and ends at the height in \eqref{eq:terminal-height}.

Every integer $t$ with $1\leq t<2m+2$ is either $t=2j$ for some
$1\leq j\leq m$, or $t=2j-1$ for some $1\leq j\leq m+1$.
The even case is covered by \eqref{eq:positive-even-height}, and the odd
case by the preceding case distinction.  Thus $h_B(t)>0$ for every
$1\leq t<2m+2$, whereas \eqref{eq:terminal-height} gives
$h_B(2m+2)=0$.  Therefore $B$ is a primitive Dyck word, and hence
$B\in\DP$.

Since the first base-$4$ digit of $p$ is $3$, one has $p\geq3$ and hence
$4p\geq12$.  Together with \eqref{eq:block-value} and $B\in\DP$, this gives
$4p\in\NpD$ and $4p\geq12$.  By the definition of $\Pset$, $p\in\Pset$.
This proves the reverse implication.
\end{proof}

The characterization gives the self-similarity used later.

\begin{corollary}\label{cor:Pclosure}
If $p\in\Pset$, then $4p+1$ and $4p+2$ also belong to $\Pset$.
\end{corollary}

\begin{proof}
By Lemma~\ref{lem:base4}, there exists $w\in\mathsf{Motz}$ such that
\[
 \operatorname{quad}(p)=3w.
\]
Then
\[
 \operatorname{quad}(4p+1)=3w1,
 \qquad
 \operatorname{quad}(4p+2)=3w2.
\]
Since
\[
 \sigma(1)=\sigma(2)=0,
\]
appending either digit preserves membership in $\mathsf{Motz}$.  The claim follows from
Lemma~\ref{lem:base4}.
\end{proof}

We next determine the smallest and largest values in each generation.

\begin{lemma}[Extremal two-coloured Motzkin values]\label{lem:extremal}
For every $\ell\geq0$ and every two-coloured Motzkin word $w$ of length $\ell$,
\begin{equation}\label{eq:motzkin-extrema}
  \frac{4^{\ell}-1}{3}\leq[w]_4\leq4^{\ell}-2^{\ell}.
\end{equation}
Both bounds are attained.  The lower bound is attained by $1^{\ell}$, and the
upper bound is attained by
\[
  w_{\max,\ell}=
  \begin{cases}
    3^a0^a,&\ell=2a,\\
    3^a2\,0^a,&\ell=2a+1.
  \end{cases}
\]
Here $1^0$ is understood to be the empty word $\lambda$.
\end{lemma}

\begin{proof}
We use induction on $\ell$.  The assertion is immediate for $\ell=0$.
Write
\[
 w=c_0c_1\cdots c_{\ell-1},
 \qquad
 c_j\in\{0,1,2,3\}.
\]

\emph{Lower bound.}
Since $\sigma(0)=-1$, the first digit cannot be $0$.  If $c_0\in\{1,2\}$, then $\sigma(c_0)=0$, so deleting the first
digit does not change any subsequent prefix weight or the total weight.
Hence the suffix $w'=c_1\cdots c_{\ell-1}$ is again a two-coloured
Motzkin word.  By the induction hypothesis for words of length $\ell-1$,
\[
  [w]_4\geq4^{\ell-1}+[w']_4
  \geq4^{\ell-1}+\frac{4^{\ell-1}-1}{3}
  =\frac{4^{\ell}-1}{3}.
\]
If $c_0=3$, then
\[
  [w]_4
  \geq[3\underbrace{00\cdots0}_{\ell-1\text{ digits}}]_4
  =3\cdot4^{\ell-1}
  \geq\frac{4^{\ell}-1}{3}.
\]
Equality is attained by $w=1^{\ell}$.

\emph{Upper bound.}
If $c_0\in\{1,2\}$, then, as above, $w'$ is again a two-coloured
Motzkin word.  By the induction hypothesis for words of length $\ell-1$,
\begin{align*}
  [w]_4
  &=[c_0c_1\cdots c_{\ell-1}]_4\\
  &\leq[2c_1\cdots c_{\ell-1}]_4\\
  &\leq2\cdot4^{\ell-1}+\bigl(4^{\ell-1}-2^{\ell-1}\bigr)\\
  &=3\cdot4^{\ell-1}-2^{\ell-1}\\
  &\leq4^{\ell}-2^{\ell}.
\end{align*}

Now suppose $c_0=3$, and let $r$ be the length of the shortest nonempty
prefix of $w$ having total weight $0$.  This prefix has the form $3u0$,
where
\[
 u=c_1\cdots c_{r-2},
\]
and let
\[
 v=c_r\cdots c_{\ell-1}
\]
be the remaining suffix.  Then
\[
 2\leq r\leq\ell,
 \qquad
 |u|=r-2,
 \qquad
 |v|=\ell-r,
\]
and $u,v\in\mathsf{Motz}$.  Thus $w=3u0v$.  By place value and the induction hypothesis,
\begin{align*}
  [w]_4
  &=3\cdot4^{\ell-1}+[u]_4\,4^{\ell-r+1}+[v]_4\\
  &\leq3\cdot4^{\ell-1}
    +\bigl(4^{r-2}-2^{r-2}\bigr)4^{\ell-r+1}
    +4^{\ell-r}-2^{\ell-r}\\
  &=4^{\ell}-2^{2\ell-r}+2^{2\ell-2r}-2^{\ell-r}\\
  &=4^{\ell}-\bigl(2^{2\ell-r}-2^{2\ell-2r}\bigr)-2^{\ell-r}.
\end{align*}
It remains to verify
\[
  2^{2\ell-r}-2^{2\ell-2r}
  \geq2^{\ell}-2^{\ell-r}.
\]
After factoring, this is
\[
  2^{2\ell-2r}(2^r-1)\geq2^{\ell-r}(2^r-1),
\]
which follows from $r\leq\ell$.  This proves the upper bound.

Finally, the displayed words $w_{\max,\ell}$ are two-coloured Motzkin words.
A geometric-series calculation gives
\[
  [w_{\max,\ell}]_4=4^{\ell}-2^{\ell},
\]
so the upper bound is attained.
\end{proof}

An element of $\Pset$ belongs to \emph{generation $d$} if its base-$4$
expansion has exactly $d$ digits.  By Lemma~\ref{lem:base4}, a
generation-$d$ element $p$ satisfies
\[
 \operatorname{quad}(p)=3w
\]
for some $w\in\mathsf{Motz}$ with $|w|=d-1$.  Therefore
\begin{equation}\label{eq:generation-value}
 p=[3w]_4=3\cdot4^{d-1}+[w]_4.
\end{equation}
The following corollary is now immediate from Lemma~\ref{lem:extremal}.

\begin{corollary}[Generation bounds]\label{cor:gen}
For every $d\geq1$, every generation-$d$ element of $\Pset$ lies in
\begin{equation}\label{eq:generation-interval}
  [g_d,U_d]
  =\left[
  3\cdot4^{d-1}+\frac{4^{d-1}-1}{3},
  4^d-2^{d-1}
  \right],
\end{equation}
and both endpoints are attained.  In particular,
\begin{enumerate}[label=\textup{(\alph*)}]
\item there is no element of $\Pset$ in $(U_d,3\cdot4^d)$;
\item for every $k\geq0$,
\begin{equation}\label{eq:three-g}
  3g_{k+1}=10\cdot4^k-1.
\end{equation}
\end{enumerate}
\end{corollary}

\begin{proof}
Let $p$ be a generation-$d$ element.  Applying
\eqref{eq:motzkin-extrema} with $\ell=d-1$ gives
\[
 \frac{4^{d-1}-1}{3}
 \leq[w]_4
 \leq4^{d-1}-2^{d-1}.
\]
Substituting these bounds into \eqref{eq:generation-value} yields
\eqref{eq:generation-interval}.  The lower
endpoint is attained by $w=1^{d-1}$, and the upper endpoint is attained by
$w=w_{\max,d-1}$.  Part~(a) follows because the next generation begins at
$g_{d+1}>3\cdot4^d$.  Finally,
\[
  3g_{k+1}
  =3\left(3\cdot4^k+\frac{4^k-1}{3}\right)
  =10\cdot4^k-1.
\]
\end{proof}


\section{Digit lifting and finite certificates}\label{sec:lifting}

The base-$4$ structure obtained in Section~\ref{sec:base4} propagates
finite interval coverings to infinite tails.  We first isolate this general
mechanism.

\begin{lemma}[Digit lifting]\label{lem:digit-lifting}
Let $k\geq1$, and let $A\subseteq\mathbb N$ satisfy
\begin{equation}\label{eq:digit-closure}
 4A+\{1,2\}\subseteq A.
\end{equation}
If $m\in kA$, then
\begin{equation}\label{eq:digit-lifting}
 [4m+k,4m+2k]\subseteq kA.
\end{equation}
\end{lemma}

\begin{proof}
Write
\[
 m=a_1+\cdots+a_k,
 \qquad a_i\in A.
\]
For every integer $r$ with $k\leq r\leq2k$, one may choose
$\eta_1,\ldots,\eta_k\in\{1,2\}$ such that
\[
 \eta_1+\cdots+\eta_k=r.
\]
By \eqref{eq:digit-closure},
\[
 4m+r
 =\sum_{i=1}^{k}(4a_i+\eta_i)
 \in kA.
\]
This proves \eqref{eq:digit-lifting}.
\end{proof}

\begin{corollary}[Propagation from a finite interval to a tail]
\label{cor:tail-propagation}
Let $k\geq3$, and let $A\subseteq\mathbb N$ satisfy
\eqref{eq:digit-closure}.  If, for some integer $M\geq1$,
\begin{equation}\label{eq:tail-seed}
 [M,4M+k-1]\subseteq kA,
\end{equation}
then
\[
 [M,\infty)\subseteq kA.
\]
\end{corollary}

\begin{proof}
We prove by strong induction on $n\geq M$ that $n\in kA$.
The assertion holds for
\[
 M\leq n\leq4M+k-1
\]
by \eqref{eq:tail-seed}.

Now let $n\geq4M+k$, and assume that every integer $j$ with
$M\leq j<n$ belongs to $kA$.  Choose the unique
\[
 \rho\in\{k,k+1,k+2,k+3\}
\]
such that $\rho\equiv n\pmod4$, and put
\[
 m:=\frac{n-\rho}{4}.
\]
Then $m$ is an integer.  Since $n\geq4M+k$ and $\rho\leq k+3$,
\[
 m\geq\frac{4M+k-(k+3)}4=M-\frac34.
\]
Hence $m\geq M$, because $m$ is an integer.  Moreover,
$\rho\geq k\geq3$, and therefore
\[
 m=\frac{n-\rho}{4}
 \leq\frac{n-k}{4}<n.
\]
The induction hypothesis gives $m\in kA$.

Finally, since $k\geq3$,
\[
 k\leq\rho\leq k+3\leq2k.
\]
Lemma~\ref{lem:digit-lifting} now yields
\[
 n=4m+\rho\in kA.
\]
The strong induction is complete.
\end{proof}

\subsection{Six-summand covering by a regular subfamily}

Define the regular base-$4$ subfamily
\[
 \mathcal E=\{0\}\cup
 \bigl\{[3\varepsilon_1\cdots\varepsilon_j]_4:
        j\geq0,\ \varepsilon_i\in\{1,2\}\bigr\},
\]
and put $\mathcal E^+=\mathcal E\setminus\{0\}$.  When $j=0$, the
suffix is the empty word, and the corresponding element is $[3]_4=3$.
Every word over $\{1,2\}$ is a two-coloured Motzkin word, so
Lemma~\ref{lem:base4} gives
\[
 \mathcal E^+\subseteq\Pset.
\]
Consequently,
\begin{equation}\label{eq:4E}
 4\cdot\mathcal E\subseteq\NpD.
\end{equation}
Appending the base-$4$ digit $1$ or $2$ also gives
\begin{equation}\label{eq:Eclosure}
 4\mathcal E^++\{1,2\}\subseteq\mathcal E^+.
\end{equation}

We use only the following finite inclusions as computational certificates.

\begin{lemma}[Finite certificate for the regular family]
\label{lem:Ecertificate}
Put
\[
\begin{aligned}
 A_0&=\{0,3,13,14\},
 &B&=\{53,54,57,58\},\\
 A&=\{3,13,14,53,54,57,58\},
 &C&=\{213,214,217,218,229,230,233,234\}.
\end{aligned}
\]
Then $A_0\subseteq\mathcal E$ and $A,B,C\subseteq\mathcal E^+$.
Moreover,
\begin{equation}\label{eq:E-small-certificate}
\begin{array}{c|c}
\text{sumset}&\text{contained interval}\\ \hline
6A_0 &[25,62]\\
B+5A_0 &[56,128]\\
2B+4A_0 &[106,172]\\
3B+3A_0 &[159,216]\\
4B+2A_0 &[212,260]
\end{array}
\end{equation}
and
\begin{equation}\label{eq:E-tail-certificate}
\begin{array}{c|c}
\text{sumset}&\text{contained interval}\\ \hline
6A &[218,260]\\
C+5A &[258,524]\\
2C+4A &[446,700]\\
3C+3A &[648,876]\\
4C+2A &[858,1052]
\end{array}
\end{equation}
hold.
\end{lemma}

\begin{proof}
Each assertion is an exhaustive finite computation of a sumset.  For a
finite set $X$, start with $S_0(X)=\{0\}$ and compute recursively
\[
 S_{i+1}(X)=S_i(X)+X.
\]
The supplementary verification script checks that every integer in each
displayed interval belongs to the corresponding sumset.
\end{proof}

\begin{proposition}\label{prop:E}
Every integer $n\geq25$ belongs to $6\mathcal E$.
\end{proposition}

\begin{proof}
The five intervals in \eqref{eq:E-small-certificate} overlap and give
\[
 [25,217]\subseteq6\mathcal E.
\]
The five intervals in \eqref{eq:E-tail-certificate} also overlap and give
\[
 [218,877]\subseteq6\mathcal E^+.
\]
By \eqref{eq:Eclosure}, the set $\mathcal E^+$ satisfies the digit-closure
condition.  Since
\[
 877=4\cdot218+6-1,
\]
Corollary~\ref{cor:tail-propagation}, applied with $k=6$ and $M=218$,
gives
\[
 [218,\infty)\subseteq6\mathcal E^+.
\]
Combining the two ranges proves the proposition.
\end{proof}

It follows from Proposition~\ref{prop:E} and \eqref{eq:4E} that every
multiple of $4$ at least $100$ is a sum of at most six primitive Dyck
numbers.

\subsection{Five-summand covering by the full halved family}

The regular family $\mathcal E$ is sufficient for the residue class $0$
modulo $4$, but the other residue class requires the full family $\Pset$.
For the finite base we use
\begin{equation}\label{eq:L-def}
 L=\{3,13,14,53,54,57,58,60\}\subseteq\Pset
\end{equation}
and
\begin{equation}\label{eq:H-def}
\begin{split}
 H=\{&213,214,217,218,220,229,230,233,234,236,\\
     &241,242,244,248\}\subseteq\Pset.
\end{split}
\end{equation}
Membership in $\Pset$ follows immediately from Lemma~\ref{lem:base4},
since the base-$4$ expansion of every displayed element has the form
$3w$ with $w\in\mathsf{Motz}$.  Enumerating the two-coloured Motzkin
words of lengths $0$, $1$, and $2$ gives exactly the eight elements of
$L$.  On the other hand, Corollary~\ref{cor:gen} gives $g_4=213$.
Hence
\begin{equation}\label{eq:P-small}
 \Pset\cap[0,211]=L.
\end{equation}

\begin{lemma}[Finite certificate for the full family]\label{lem:finitefive}
The following inclusions hold:
\begin{equation}\label{eq:finite-interval-certificate}
\begin{array}{c|c}
\text{sumset}&\text{contained interval}\\ \hline
5L &[215,254]\\
H+4L &[245,486]\\
2H+3L &[439,674]\\
3H+2L &[645,862]\\
4H+L &[855,1046].
\end{array}
\end{equation}
Moreover,
\begin{equation}\label{eq:finite-five-lower-certificate}
 209,210,211\notin F_5(L).
\end{equation}
\end{lemma}

\begin{proof}
Equation~\eqref{eq:finite-interval-certificate} is a finite-sumset
inclusion certificate, while \eqref{eq:finite-five-lower-certificate}
is an exhaustive computation of $F_5(L)$ in the relevant finite range.
Both are verified by the same recursive finite-set addition used in
Lemma~\ref{lem:Ecertificate} and are reproduced by the supplementary
verification script.
\end{proof}

The five intervals in Lemma~\ref{lem:finitefive} overlap, and hence
\begin{equation}\label{eq:fiveP-seed}
 [215,1046]\subseteq5\Pset.
\end{equation}

\begin{proposition}\label{prop:fiveP}
Every integer $n\geq215$ is a sum of exactly five elements of $\Pset$.
Consequently,
\[
 [212,\infty)\subseteq F_5(\Pset).
\]
The lower endpoint is sharp: none of $209,210,211$ belongs to
$F_5(\Pset)$.
\end{proposition}

\begin{proof}
By Corollary~\ref{cor:Pclosure},
\[
 4\Pset+\{1,2\}\subseteq\Pset.
\]
Equation~\eqref{eq:fiveP-seed} contains, in particular,
\[
 [215,864]\subseteq5\Pset.
\]
Since
\[
 864=4\cdot215+5-1,
\]
Corollary~\ref{cor:tail-propagation}, applied with $k=5$ and $M=215$,
gives
\[
 [215,\infty)\subseteq5\Pset.
\]
Moreover,
\[
\begin{aligned}
212&=53+53+53+53,\\
213&=53+53+53+54,\\
214&=53+53+54+54,
\end{aligned}
\]
so $[212,\infty)\subseteq F_5(\Pset)$.  Finally,
\eqref{eq:P-small} and \eqref{eq:finite-five-lower-certificate} give
\[
 209,210,211\notin F_5(\Pset).
\]
\end{proof}

\section{The complete exceptional set}\label{sec:exceptions}

The decomposition
\[
 \NpD=\{0,2\}\cup4\Pset
\]
converts summand-count conditions uniformly into finite-sum conditions on
$\Pset$.

\begin{lemma}[Mod-$4$ summand criterion]\label{lem:mod4-sumcriterion}
Let $\epsilon\in\{0,1\}$, let $h\geq0$, and put
$N=4m+2\epsilon$.  Then $N$ is a sum of at most $h$ positive primitive
Dyck numbers if and only if
\begin{equation}\label{eq:mod4-sumcriterion}
 m\in
 \bigcup_{\substack{0\leq s\leq h\\
                    s\equiv\epsilon\pmod2}}
 \left(
  \frac{s-\epsilon}{2}+F_{h-s}(\Pset)
 \right).
\end{equation}
\end{lemma}

\begin{proof}
Let $s$ be the number of summands equal to $2$.  By
Lemma~\ref{lem:mod4}, every remaining positive summand has the form
$4p_i$ with $p_i\in\Pset$.  Reduction modulo $4$ gives
$s\equiv\epsilon\pmod2$, and
\[
 4m+2\epsilon=2s+4\sum_i p_i
\]
is equivalent to
\[
 m=\frac{s-\epsilon}{2}+\sum_i p_i.
\]
There are at most $h-s$ remaining summands, which proves necessity.
Conversely, any representation belonging to one of the sets in
\eqref{eq:mod4-sumcriterion} gives a representation of $N$ after
replacing each $p_i$ by $4p_i$ and adjoining $s$ copies of $2$.
\end{proof}

\begin{corollary}\label{cor:sixcriterion}
Let $N=4m+2$.  Then $N$ is a sum of at most six positive primitive Dyck
numbers if and only if
\begin{equation}\label{eq:sixcriterion}
 m\in F_5(\Pset)
 \cup\bigl(1+F_3(\Pset)\bigr)
 \cup\bigl(2+F_1(\Pset)\bigr).
\end{equation}
\end{corollary}

\begin{proof}
Apply Lemma~\ref{lem:mod4-sumcriterion} with
$\epsilon=1$ and $h=6$, and use $s=1,3,5$.
\end{proof}

\begin{lemma}[Small multiples of four]\label{lem:smallmultiples}
Among the positive multiples of $4$ below $100$, the only integer not
representable by at most six positive primitive Dyck numbers is $44$.
Moreover,
\[
 44=12+12+12+2+2+2+2.
\]
\end{lemma}

\begin{proof}
The positive primitive Dyck numbers below $100$ are precisely
\[
 2,12,52,56.
\]
An exhaustive computation of sums of at most six elements of this finite
set, truncated below $100$, shows that the only missing positive multiple
of $4$ is $44$.  The displayed identity gives a seven-summand
representation.  This finite check is also included in the supplementary
verification script.
\end{proof}

\begin{lemma}[Finite exceptional-set certificate]\label{lem:finiteexception}
Let $L$ be the set defined in \eqref{eq:L-def}.  Then
\begin{equation}\label{eq:finiteexception}
\begin{split}
 [0,211]\setminus\Bigl(
 F_5(L)&\cup(1+F_3(L))\cup(2+F_1(L))\Bigr)\\
 &=\{8,11,24,38,49,50,51,209,210,211\}.
\end{split}
\end{equation}
\end{lemma}

\begin{proof}
Starting from the finite set $L$, compute $F_1(L)$, $F_3(L)$, and
$F_5(L)$ exactly by recursive finite-set addition, truncate the result to
$[0,211]$, and take the union and its complement.  This gives
\eqref{eq:finiteexception}; the computation is reproduced by the
supplementary verification script.
\end{proof}

\begin{proof}[Proof of Theorem~\ref{thm:classification}]
Suppose first that $N\equiv0\pmod4$.  If $N\geq100$, write $N=4n$.
Then $n\geq25$, so Proposition~\ref{prop:E} and \eqref{eq:4E} give a
representation of $N$ with at most six primitive Dyck summands.
Lemma~\ref{lem:smallmultiples} treats the remaining positive multiples
of $4$ and shows that $r(44)=7$.

Now let $N\equiv2\pmod4$, say $N=4m+2$.  If $N\geq850$, then
$m\geq212$, so Proposition~\ref{prop:fiveP} gives
$m\in F_5(\Pset)$.  Corollary~\ref{cor:sixcriterion} gives a
representation of $N$ with at most six primitive Dyck summands.  Together
with the preceding paragraph, this already shows that every even
$N\geq848$ has such a representation.

It remains to classify the integers $N\leq846$ with
$N\equiv2\pmod4$.  Here $0\leq m\leq211$.  By
\eqref{eq:P-small}, Lemma~\ref{lem:finiteexception}, and
Corollary~\ref{cor:sixcriterion}, the corresponding values $N=4m+2$
that are not sums of at most six primitive Dyck numbers are precisely
\[
 34,46,98,154,198,202,206,838,842,846.
\]
Thus these ten numbers, together with $44$, are exactly the positive even
integers that are not sums of at most six primitive Dyck numbers.

All of them except $46$ have seven-summand representations:
\[
\begin{aligned}
34&=12+12+2+2+2+2+2,\\
44&=12+12+12+2+2+2+2,\\
98&=56+12+12+12+2+2+2,\\
154&=52+52+12+12+12+12+2,\\
198&=56+52+52+12+12+12+2,\\
202&=56+56+52+12+12+12+2,\\
206&=56+56+56+12+12+12+2,\\
838&=240+240+240+52+52+12+2,\\
842&=240+240+240+56+52+12+2,\\
846&=240+240+240+56+56+12+2.
\end{aligned}
\]
Proposition~\ref{prop:46} shows that $46$ requires exactly eight
summands.  This proves all assertions of the theorem.
\end{proof}

\begin{corollary}\label{cor:threshold}
The integer $848$ is the least even integer $T$ such that every even
integer $N\geq T$ is a sum of at most six primitive Dyck numbers.
\end{corollary}

\begin{proof}
Theorem~\ref{thm:classification} shows that six summands suffice for every
even integer at least $848$, whereas $846$ requires seven summands.
\end{proof}

\begin{corollary}\label{cor:eight}
Every nonnegative even integer is a sum of at most eight elements of
$\NpD$, and $46$ is the unique even integer that is not a sum of at most
seven elements of $\NpD$.
\end{corollary}

\begin{proof}
The assertion for positive integers follows immediately from
Theorem~\ref{thm:classification}; the integer $0$ is the empty sum.
\end{proof}

For $A\subseteq2\mathbb N_0$, we call $A$ a \emph{relative additive basis of
order at most $h$ for $2\mathbb N_0$} if every element of $2\mathbb N_0$ is a
sum of at most $h$ elements of $A$.  Its relative order is \emph{exactly $h$}
if $h$ is the least integer with this property.  We similarly call $A$ a
\emph{relative asymptotic additive basis of order at most $h$} if every
sufficiently large element of $2\mathbb N_0$ is a sum of at most $h$ elements
of $A$.

Theorem~\ref{thm:classification} gives the complete exceptional set for
six-summand representability.  Corollary~\ref{cor:eight} shows that $\NpD$ is
a relative additive basis of exact order eight for $2\mathbb N_0$, while
Corollary~\ref{cor:threshold} identifies the sharp six-summand threshold.
The next section explains the difference between the two residue classes and
proves that the relative asymptotic order is exactly six.

\section{Recurring gaps and asymptotic optimality}\label{sec:asymptotic}

For multiples of $4$, the five-summand representations in $\Pset$ may be
used directly.  For integers congruent to $2$ modulo $4$, however, an odd
number of copies of the exceptional summand $2$ must be used.  This
distinction raises the relative asymptotic order of the full primitive
family from $5$ to $6$.

\begin{corollary}[Five-summand criterion]\label{cor:fivecrit}
Let $N=4m+2$ with $m\geq3$.  Then $r(N)\leq5$ if and only if
\begin{equation}\label{eq:fivecrit}
 m\in F_4(\Pset)\cup\bigl(1+F_2(\Pset)\bigr).
\end{equation}
\end{corollary}

\begin{proof}
Apply Lemma~\ref{lem:mod4-sumcriterion} with $\epsilon=1$ and $h=5$.
The values $s=1,3,5$ give
\[
 F_4(\Pset)\cup\bigl(1+F_2(\Pset)\bigr)\cup\{2\}.
\]
The last set is irrelevant because $m\geq3$.
\end{proof}

The following lemma isolates the lower-bound mechanism created by the
gaps between generations.

\begin{lemma}[Recurring obstruction]\label{lem:obstruction}
For every integer $k\geq2$, put
\[
 x_k=10\cdot4^k-2.
\]
Then
\begin{equation}\label{eq:obstruction}
 x_k\notin F_4(\Pset)
 \qquad\text{and}\qquad
 x_k-1\notin F_2(\Pset).
\end{equation}
\end{lemma}

\begin{proof}
Fix $k\geq2$, and abbreviate $x=x_k$.  The least element of generation
$k+2$ is
\[
 g_{k+2}>3\cdot4^{k+1}=12\cdot4^k>x.
\]
Consequently, every element of $\Pset$ that is at most $x$ has generation
at most $k+1$.  By Corollary~\ref{cor:gen}, every such element is either
\emph{small}, of generation at most $k$ and hence at most
\[
 U_k=4^k-2^{k-1},
\]
or \emph{large}, of generation $k+1$ and hence between $g_{k+1}$ and
$U_{k+1}=4^{k+1}-2^k$.

First consider $x-1=10\cdot4^k-3$.  It is larger than $U_{k+1}$ and
smaller than $g_{k+2}$, so it does not itself belong to $\Pset$.
Moreover, any sum of two eligible elements is at most
\[
 2U_{k+1}=8\cdot4^k-2^{k+1}<10\cdot4^k-3=x-1.
\]
Thus $x-1\notin F_2(\Pset)$.

Now suppose, for contradiction, that
\[
 x=q_1+\cdots+q_t,
 \qquad t\leq4,
 \qquad q_i\in\Pset.
\]
Let $N_{\rm b}$ be the number of large summands.  If $N_{\rm b}\leq2$,
then filling the remaining positions with the largest possible small
summands gives
\begin{align*}
 x
 &\leq2U_{k+1}+2U_k\\
 &=2(4^{k+1}-2^k)+2(4^k-2^{k-1})\\
 &=10\cdot4^k-3\cdot2^k
 <10\cdot4^k-2=x,
\end{align*}
a contradiction.  Hence $N_{\rm b}\geq3$.  But then the three large
summands alone have sum at least
\[
 3g_{k+1}=10\cdot4^k-1=x+1
\]
by \eqref{eq:three-g}, again a contradiction.  Therefore
$x\notin F_4(\Pset)$.
\end{proof}

\begin{corollary}[Order of $\Pset$]\label{cor:orderfive}
The set $\Pset$ is an asymptotic additive basis of order exactly $5$ for
the nonnegative integers.  More precisely, every integer $n\geq212$ lies
in $F_5(\Pset)$, whereas $F_4(\Pset)$ and $F_3(\Pset)$ each omit
infinitely many integers.
\end{corollary}

\begin{proof}
Proposition~\ref{prop:fiveP} gives
$[212,\infty)\subseteq F_5(\Pset)$.  For every $k\geq2$,
Lemma~\ref{lem:obstruction} gives $x_k\notin F_4(\Pset)$.  Since
$x_k\to\infty$, four summands do not suffice asymptotically.  The same is
true for three summands because $F_3(\Pset)\subseteq F_4(\Pset)$.
\end{proof}

\begin{corollary}[Multiples of $4$ requiring five summands]
\label{cor:five-multiples}
For every integer $k\geq3$,
\[
 M_k=10\cdot4^{k+1}-8
\]
satisfies $r(M_k)=5$.  Hence infinitely many multiples of $4$ require
exactly five positive primitive Dyck summands.
\end{corollary}

\begin{proof}
Write $M_k=4x_k$.  Since $k\geq3$, one has $x_k\geq212$, and
Proposition~\ref{prop:fiveP} gives $x_k\in F_5(\Pset)$.  Multiplying such
a representation by $4$ shows that $r(M_k)\leq5$.

Suppose that $r(M_k)\leq4$.  Applying
Lemma~\ref{lem:mod4-sumcriterion} with $\epsilon=0$ and $h=4$ gives
\[
 x_k\in
 F_4(\Pset)
 \cup\bigl(1+F_2(\Pset)\bigr)
 \cup\{2\}.
\]
The first two alternatives contradict Lemma~\ref{lem:obstruction}, and
the last is impossible because $x_k>2$.  Thus $r(M_k)\geq5$, and equality
follows.
\end{proof}

\begin{proof}[Proof of Theorem~\ref{thm:asymptotic-order}]
Fix $k\geq2$ and put $m=x_k=10\cdot4^k-2$.  Then
\[
 N_k=4m+2=10\cdot4^{k+1}-6.
\]
By Lemma~\ref{lem:obstruction},
\[
 m\notin F_4(\Pset),
 \qquad
 m-1\notin F_2(\Pset).
\]
Corollary~\ref{cor:fivecrit} therefore gives $r(N_k)\geq6$.  The integer
$N_k$ is not one of the eleven exceptions in
Theorem~\ref{thm:classification}, so $r(N_k)\leq6$.  Hence $r(N_k)=6$.
Since $N_k\to\infty$, no $h\leq5$ bounds $r(N)$ for all sufficiently
large even $N$, while Theorem~\ref{thm:classification} gives the upper
bound $h=6$.  Therefore $h_{\mathrm{asym}}=6$.
\end{proof}

\begin{remark}[The two residue classes]
If $N=4n$, then $n\in F_5(\Pset)$ for all sufficiently large $n$, so
$r(N)\leq5$ eventually; Corollary~\ref{cor:five-multiples} shows that
equality occurs infinitely often.  If $N=4m+2$, the number of copies of
$2$ must be odd.  For $m=x_k$, the one-copy case would require
$m\in F_4(\Pset)$ and the three-copy case would require
$m-1\in F_2(\Pset)$; Lemma~\ref{lem:obstruction} excludes both.  This
residue class raises the relative asymptotic order from $5$ to $6$.
\end{remark}

\section{OEIS sequences and computational verification}
\label{sec:oeis}

Three OEIS sequences occur naturally.  The totally balanced numbers form
\seqnum{A014486}~\cite{OEIS-A014486}; the positive primitive Dyck numbers
form \seqnum{A057547}~\cite{OEIS-A057547}; and the central sequence of
this paper is
\[
 a(n)=r(2n),
\]
namely \seqnum{A395858}~\cite{OEIS-A395858}.  Theorem~\ref{thm:classification}
determines all terms exceeding six:
\[
 a(n)=8\iff n=23,
\]
and
\[
 a(n)=7\iff
 n\in\{17,22,49,77,99,101,103,419,421,423\}.
\]
Theorem~\ref{thm:asymptotic-order} supplies an explicit infinite
subsequence at level six:
\begin{equation}\label{eq:A395858-six-subsequence}
 a(5\cdot4^{k+1}-3)=6\qquad(k\geq2).
\end{equation}

The program
\href{https://github.com/growupkuriyama-hub/lean_cfg_project/blob/3525cfb/LeanCfgProject/PrimDyck/A395858.py}
{\path{A395858.py} at commit \texttt{3525cfb}}
reproduces the displayed initial terms and generates the first $50{,}000$
terms of A395858.  The finite certificates in this paper are reproduced by
the supplementary script
\href{https://github.com/growupkuriyama-hub/lean_cfg_project/blob/f38e719/LeanCfgProject/PrimDyck/verify_certificates_for_paper.py}
{\path{verify_certificates_for_paper.py}}.
The filename is unchanged from the preceding version.  The v6-compatible
script uses only exact integer arithmetic and the Python standard library.
It verifies the regular-family certificates
\eqref{eq:E-small-certificate}--\eqref{eq:E-tail-certificate}, the
full-family certificate \eqref{eq:finite-interval-certificate}, the sharp
lower certificate \eqref{eq:finite-five-lower-certificate}, the small
multiples of $4$, and the finite exceptional-set identity
\eqref{eq:finiteexception}.  It also performs an independent
dynamic-programming cross-check of $r(N)$ up to a user-selected bound.
The version at commit \texttt{f38e719} passed Python CI run \#37.  All
infinite-range assertions are proved theoretically in the text by the
digit-lifting lemma, the generation bounds, and the recurring-obstruction
argument.

\section{Concluding remarks}

The results determine three exact orders associated with primitive Dyck
numbers.  The relative additive order on all nonnegative even integers is
$8$, the relative asymptotic order is $6$, and the asymptotic additive
order of $\Pset$ is $5$.  The same base-$4$ self-similarity that propagates
finite sumset certificates to infinite tails also produces the recurring
obstruction $x_k=10\cdot4^k-2$.

It remains to characterize the complete set of even integers $N$ for which
$r(N)=6$.  Computations suggest additional block structure near the points
$10\cdot4^k$, but the exact widths and endpoints of those blocks are not
yet proved.  A broader question is whether analogous regular-subfamily
upper bounds and generation-gap lower bounds occur for other indecomposable
digital languages.

\begin{thebibliography}{9}

\bibitem{BHS}
J.~P. Bell, K.~G. Hare, and J.~Shallit,
When is an automatic set an additive basis?,
\emph{Proc. Amer. Math. Soc. Ser. B} \textbf{5} (2018), 50--63.

\bibitem{BLS}
J.~P. Bell, T.~F. Lidbetter, and J.~Shallit,
Additive number theory via approximation by regular languages,
\emph{Internat. J. Found. Comput. Sci.} \textbf{31} (2020), 667--687.

\bibitem{MNRS}
P.~Madhusudan, D.~Nowotka, A.~Rajasekaran, and J.~Shallit,
Lagrange's theorem for binary squares,
in \emph{43rd International Symposium on Mathematical Foundations of
Computer Science}, Leibniz Int. Proc. Inform., Vol.~117, Schloss
Dagstuhl--Leibniz-Zentrum f{\"u}r Informatik, 2018, pp.~18:1--18:14.

\bibitem{Nathanson}
M.~B. Nathanson,
\emph{Additive Number Theory: The Classical Bases},
Grad. Texts in Math., Vol.~164, Springer, 1996.

\bibitem{OEIS-A014486}
OEIS Foundation Inc.,
\emph{The On-Line Encyclopedia of Integer Sequences, A014486},
\url{https://oeis.org/A014486}.

\bibitem{OEIS-A057547}
OEIS Foundation Inc.,
\emph{The On-Line Encyclopedia of Integer Sequences, A057547},
\url{https://oeis.org/A057547}.

\bibitem{OEIS-A395858}
OEIS Foundation Inc.,
\emph{The On-Line Encyclopedia of Integer Sequences, A395858},
\url{https://oeis.org/A395858}.

\bibitem{RSS}
A.~Rajasekaran, J.~Shallit, and T.~Smith,
Sums of palindromes: an approach via automata,
in \emph{35th Symposium on Theoretical Aspects of Computer Science},
Leibniz Int. Proc. Inform., Vol.~96, Schloss Dagstuhl--Leibniz-Zentrum
f{\"u}r Informatik, 2018, pp.~54:1--54:12.

\end{thebibliography}


\begin{flushleft}
Takayuki Kuriyama\\
Growup Learning Academy\\
4-8-3 Shakujii-machi\\
Nerima-ku, Tokyo 177-0041\\
Japan\\
\textit{Email address}: \href{mailto:growup.kuriyama@gmail.com}
{\texttt{growup.kuriyama@gmail.com}}
\end{flushleft}

\end{document}
